\documentclass[11pt]{article}
\usepackage[margin=0.9in]{geometry}
\usepackage{hyperref}
\usepackage{enumitem}
\usepackage{xcolor}
\usepackage{amssymb}
\usepackage{fancyhdr}

\setlist[itemize]{leftmargin=*,topsep=2pt,itemsep=1pt}
\setlist[enumerate]{leftmargin=*,topsep=2pt,itemsep=1pt}
\pagestyle{fancy}
\fancyhf{}
\rhead{ScuffedRDMA questions}
\lhead{Nathan Gopee, SUNY New Paltz}
\rfoot{\thepage}

\title{Questions: RDMA transport for disaggregated LLM inference\\[4pt]
       \large Discussion guide for team meeting}
\author{Nathan Gopee\\MS Computer Science, SUNY New Paltz\\\texttt{gopeen1@newpaltz.edu}}
\date{\today}

\begin{document}
\maketitle
\thispagestyle{fancy}

\hrule
\vspace{6pt}

\section*{Notes to myself before the meeting}

\textbf{What to cover in the first three minutes (before I start asking
things).}

\begin{enumerate}
\item \textbf{Thesis in one sentence.} An adaptive RDMA middleware for
disaggregated LLM inference over commodity 10\,GbE Ethernet, targeting
the KV-cache-transfer term of Time-to-First-Token in vLLM
prefill/decode serving.

\item \textbf{Why the problem is real.} On a single RC queue pair, an
8\,$\mu$s metadata message and a 16\,MB KV block share one work queue,
so the small transfer stalls behind the big one --- head-of-line
blocking in software. InfiniBand solves it with hardware QoS. I'm
building a software path over SoftRoCE that splits the queue pair into
two priority classes.

\item \textbf{Four components I've built.}
  \begin{itemize}
  \item \textbf{Dual QP pool} --- one protection domain, two RC QPs
  (hot: busy-poll CQ; cold: sleep-poll).
  \item \textbf{WFA classifier} --- Work-First-Adaptive policy that
  labels each transfer hot or cold from tensor size plus the
  prefill-vs-decode phase signal.
  \item \textbf{PMP controller} --- Pontryagin's Maximum Principle
  gives a bang-bang solution for bandwidth allocation between two
  queues. Switching function
  $S = \alpha q_H C \mu_H - \beta q_C C \mu_C$ uses queue depth and
  measured service rate.
  \item \textbf{scuffedQuant} --- two-stage KV compression
  (PolarQuant + QJL sign-sketch) at 3 bits with no calibration.
  91\,\% top-8 attention overlap on Granite 3.3-2B KV cache.
  \end{itemize}

\item \textbf{Testbed (consumer-grade).}
  \begin{itemize}
  \item Chimera: 3$\times$~RTX~3090, Aquantia AQC107 10\,GbE, SoftRoCE
  (\texttt{rxe0}).
  \item Cerberus: 2$\times$~RTX~5090, Intel X710 10\,GbE, SoftRoCE +
  iWARP SR-IOV.
  \item Measured: 0.92 Gb/s SoftRoCE loopback, 12.6\,$\mu$s single-QP
  p50, +0.6\,$\mu$s dual-QP overhead, 8\,$\mu$s cross-node decode
  latency.
  \end{itemize}

\item \textbf{Upstream contributions.} Six pull requests to
\texttt{openucx/ucx} (\#11304--\#11309), all reviewed in Update~5.
SoftRoCE wireup fix (\#11307) is the highest-impact one; it removes
UCX's hard dependency on UD for the wireup auxiliary transport.

\item \textbf{What I want from this meeting.} Not validation --- a
reality check from people with more production miles than me. Where is
my model wrong? What am I measuring that I shouldn't trust? What
production failure modes am I about to walk into?
\end{enumerate}

\textbf{Don't do.}
\begin{itemize}
\item Don't dump the full PDF stack on intro. Send Update 4 + Update 5
at most.
\item Don't lead with the compression numbers. Storage engineers have
seen a thousand compressors. Lead with the classify-then-route
transport story.
\item Don't over-claim. Every number I quote has a JSON file behind it
under \texttt{benchmarks/results/}. If I don't have the JSON, I say
"that's a target, not a measurement."
\end{itemize}

\vspace{8pt}
\hrule
\vspace{6pt}

\section{RDMA fundamentals I want to sanity-check}

\begin{enumerate}
\item My dual-QP design treats head-of-line blocking at the work-queue
level. Am I conflating two problems --- host-side queuing and
wire-side congestion? Where does PFC / ECN / DCQCN actually matter for
my traffic mix (many small metadata ops + sparse large KV blocks)?

\item RoCEv2 on a lossy 10\,GbE link without PFC: what symptoms should
I expect before I see throughput collapse? Drops vs. retransmits vs.
RNR NAKs --- which one do I watch first?

\item Direct-connect (no switch) vs. switched RoCE: what does the
failure mode look like when someone plugs a switch in later? Anything
beyond the ARP/AH story the r/HPC thread already covered?

\item SoftRoCE (\texttt{rdma\_rxe}) as a stand-in for hardware RoCE ---
which of my measurements are load-bearing on real hardware, and which
ones are rxe artifacts? I've hit the GRH init issue already; what
else?

\item iWARP SR-IOV on the 5090 side is theoretical for me right now.
Is iWARP still worth the effort for LLM inference traffic patterns, or
has everyone effectively standardized on RoCEv2?

\item GID index selection: libmesh-rdma (the project I ported from)
had to scan for the IPv4-mapped GID because hard-coding index 0 was
fragile. Is the right answer still to scan, or is there a stable
ioctl/sysfs path production sites use?

\item MTU: I'm running the rxe default. Should I push to 4096 or 9000
(jumbo) on consumer NICs? What breaks first when I do?

\item RC vs. RD vs. UD for KV traffic: I picked RC because PagedAttention
block transfers are point-to-point and reliable. Am I leaving
performance on the table by not exploring UD for the hot-path control
messages specifically?

\item Queue-pair lifetime: I create QPs at connection setup and hold
them open. At a few hundred endpoints per host is that OK, or do I
need a QP pool with per-peer aging like Lustre's LNET peer table?

\item One-sided WRITE vs. SEND/RECV for KV blocks: WRITE wins on the
latency side but requires the sender to know the receiver's rkey and
address. At what fan-out does the bookkeeping overhead for WRITE start
to beat the simpler SEND/RECV path?
\end{enumerate}

\section{pyverbs / libibverbs / rdma-core}

\begin{enumerate}
\item Is pyverbs production-viable for the hot path, or am I fooling
myself because SoftRoCE is hiding the GIL cost? At what message rate
does Python lose to a C/Rust port?

\item Memory-region registration: I register a large slab once, then
issue WRITEs into offsets. Is that the right default, or should I be
doing on-demand registration + MR cache per vLLM block?

\item IOMMU interaction with pinned MRs on consumer motherboards: any
gotchas with Ryzen / consumer Intel? I haven't had issues yet but I
haven't run long-duration tests.

\item Completion queue sizing: my hot pool has a CQ depth of 128
against a send-queue depth of 64. Is that overbuilt, under-built, or
irrelevant at my scale?

\item Busy-poll vs. event-driven CQ: hot pool busy-polls (one core
pinned). On a 3-3090 box that burns a core continuously. Is that
acceptable on a production AI box, or is there an adaptive-poll
pattern I should be using?

\item Shared receive queues (SRQ): do they matter for my topology
(two nodes, fan-in = 1)? Or is it only at larger fan-in that SRQ
starts paying off?

\item \texttt{ibv\_post\_send} batching: I submit one WR at a time.
Should I be batching (posting a linked list) for the cold pool?

\item \texttt{ibv\_poll\_cq} batch size: I poll up to 16 at a time.
Is there a number I should justify rather than guess?

\item Extended verbs (\texttt{ibv\_exp\_*}) vs. the ``standard'' verbs
API --- am I missing anything important by staying on the standard
API?

\item rdma-core vs. MLNX OFED: consumer ConnectX-4 in Tower 2 is
coming. Is upstream rdma-core enough or do I need the Mellanox OFED
kernel modules for anything specific (GDR, GID assignment, counters)?
\end{enumerate}

\section{Block sizes, message granularity, KV transport}

\begin{enumerate}
\item vLLM's PagedAttention uses 16-token blocks. At \texttt{head\_dim}
= 128, \texttt{n\_kv\_heads} = 32, that's $16 \cdot 2 \cdot 32 \cdot
128 \cdot 2 = 256$\,KB per block (FP16). Is 256\,KB the natural
unit for a single RDMA WRITE, or should I be aggregating blocks?

\item Scatter/gather: one WRITE with an SGE list covering $N$ blocks,
or $N$ WRITEs? Where's the crossover on a 10\,GbE link?

\item RDMA Read vs. Write for KV fetch: decode is the puller. Is it
better for decode to issue READs (pull model, self-paced) or for
prefill to issue WRITEs (push model, sender-paced)?

\item Message inline vs. non-inline: my hot-path control messages are
small (tens of bytes). At what size does inline stop helping and start
hurting?

\item Unsignaled completions: I signal every completion. Is that
wasting CQ cycles I should be reclaiming via a signaled-every-N
strategy?

\item Alignment: PagedAttention blocks aren't necessarily aligned to
NIC PCIe burst boundaries. Does that show up as measurable overhead
in practice, or is it noise?

\item FlashInfer's MXFP4 KV cache layout: does that change the RDMA
picture at all, or is it opaque to transport?

\item scuffedQuant compresses to 3 bits, so a 256\,KB block becomes
$\approx 32$\,KB. Does dropping into a different size class for the
WFA classifier change anything I should know about?

\item On-GPU vs. host-memory staging: I stage through host memory
today (no GDR). How much are my latency numbers inflated by the
PCIe round-trip per block?

\item Maximum outstanding WRs per QP: what's a sensible value to set
on the hot QP so I don't starve the RC pipeline but also don't build
a deep queue that defeats the whole hot/cold split?
\end{enumerate}

\section{vLLM integration}

\begin{enumerate}
\item The \texttt{KvConnector} interface: how stable is it really?
What's churning upstream that will break my integration if I pin too
loose?

\item Disaggregated prefill/decode in vLLM: is the expected production
path through NIXL (NVIDIA Inference Xfer Library), through llm-d, or
through a direct \texttt{KvConnector} subclass? Am I picking the right
integration point?

\item PagedAttention block table translation across nodes: prefill and
decode have different physical block IDs. Is there a canonical way to
exchange the mapping, or is everyone rolling their own?

\item Tensor parallelism across a heterogeneous pair (3$\times$3090 vs.
2$\times$5090): does that even work end-to-end in vLLM, or do I need
pipeline parallelism to bridge the mismatch?

\item MXFP4 on gpt-oss-120b: I have a 104.4 tok/s TCP baseline on
Chimera. What should my RDMA target actually be before I call it a
result?

\item KV cache reuse across requests: with LMCache, cross-request
reuse produces a very different traffic pattern (bulk read of prior
context). Does my hot/cold split hold for that, or does it degenerate
to all-cold?

\item What's the lowest-friction way to drop libscuffedrdma in front
of an unmodified vLLM for a first-pass measurement? A shim
\texttt{TorchDistributed} backend? A \texttt{NIXL} plug-in?

\item Which upstream vLLM PRs should I be tracking (not just
referencing) for disaggregated prefill/decode?
\end{enumerate}

\section{I/O and RDMA testing methodology}

\begin{enumerate}
\item \texttt{ib\_write\_bw}, \texttt{ib\_write\_lat},
\texttt{ib\_send\_bw}: which of these actually correlate with
my application-level numbers, and which are misleading me?

\item \texttt{perftest} vs. application-level microbenchmarks: where
is each one strong, and what does each one hide?

\item Multi-threaded stress tests: I don't have one yet. What should
it look like? Concurrent WRITEs from N threads into one QP pool, or
N QP pools?

\item Long-duration stability: 24\,h of continuous inference. What
counters should I be watching? \texttt{perfquery}? \texttt{nvidia-smi
dmon}? Port-level error rates?

\item Latency percentiles: I report p50/p95/p99. Is p99.9 the right
tail to care about for inference dispatch, or is p99 enough?

\item Microbenchmark noise floor: how do I distinguish a real +0.6
$\mu$s dual-QP overhead from measurement noise at the sub-microsecond
scale?

\item Fault injection: should I be deliberately triggering QP errors
(receiver posts no buffers, sender blasts anyway) and measuring
recovery time? If so, what's the accepted way to do that?

\item Cold boot vs. warm run: I pre-warm the pools before measuring.
Is there a canonical warm-up protocol I should follow so my numbers
are comparable with published UCX/NCCL numbers?

\item CPU pinning, NUMA, IRQ affinity: I set core pinning for the
busy-poll thread but I haven't tuned NUMA. What's the real-world
impact?

\item Profiling tools: I use \texttt{perf} + \texttt{py-spy}.
Is there an RDMA-specific profiling toolkit I should know (MLNX
NeoHost? UCX profiling?)?

\item Regression testing: how do I prevent a future me from shipping
a latency regression? What's the standard CI pattern for RDMA code?
\end{enumerate}

\section{Production system lessons (Lustre, GPFS, UCX at scale)}

\begin{enumerate}
\item Lustre's LNET credit accounting between small-message and
bulk-message queues --- is that isomorphic to my PMP bang-bang
controller, or is there something LNET does that I'm missing?

\item Spectrum Scale NSD client uses RDMA. Does it run multiple QPs
per NSD pair (priority classes), or a single QP with in-band
prioritization?

\item At what cluster size does single-QP head-of-line blocking
become provably painful in production? 10 nodes? 100? Or is it the
traffic mix, not the node count, that matters?

\item If my dual-QP + WFA + PMP abstraction isn't novel in the
filesystem world, where should I be looking for prior art? I'd rather
cite it than rediscover it.

\item UCX upstream: the six PRs I submitted --- are those the right
six, or are the important issues somewhere else in the codebase I
haven't read?

\item What's the right failure model for an RDMA transport in an
inference pipeline? An NSD client can retry for minutes; a decode step
has a tens-of-milliseconds budget. What do production systems do at
that gap?

\item Network tuning for mixed traffic: production Lustre/GPFS
clusters segregate metadata traffic onto its own IB partition. Is
that the filesystem analog of my hot QP, or are the two mechanisms
solving different problems?

\item GPUDirect RDMA on consumer GPUs: I have ConnectX-4 on Tower 2
and V100s. Is it worth the integration pain on consumer motherboards
without BlueField, or should I stage through host RAM indefinitely?
\end{enumerate}

\section{Which open issues should I dive into the most?}

The repository has 15 open issues. I have an ordering, but I'd rather
hear where experienced engineers would start.

\textbf{My ranking (highest return first):}

\begin{enumerate}
\item \textbf{\#35: TTFT reduction measurement.} The only measurement
that matters for the thesis claim. Currently blocked by \#34.

\item \textbf{\#34: vLLM disaggregated prefill/decode with RDMA.}
The integration that makes \#35 possible. Highest implementation risk.

\item \textbf{\#57: Multi-threaded concurrent workload benchmark.}
Demonstrates dual-QP's value. Without this, Update~4's claim about
HOL blocking is theoretical.

\item \textbf{\#19: GPT-oss-120b RDMA vs. TCP benchmark.} First
real ``RDMA beats TCP'' data point. Needs Tower~1 + Tower~2 with
ConnectX-4 working.

\item \textbf{\#37: 24-hour RoCE stability test.} Surfaces real
failure modes a microbenchmark misses.

\item \textbf{\#60: WFA trace-replay test.} Gives me a reproducible
microbench for classifier regressions. Low implementation cost, high
long-term payoff.

\item \textbf{\#59: PMP stability sweep.} Answers ``is the bang-bang
safe to ship?'' Must do before \#38 (ablation).

\item \textbf{\#58: scuffedQuant extended validation.} Widens the
model coverage beyond Granite 3.3-2B, strengthens the compression
claim.

\item \textbf{\#38: Full ablation study.} Thesis Chapter~5's results
table. Blocked by \#35 + \#57.

\item \textbf{Thesis chapters (\#39--\#44).} Writing cadence, no
technical risk.
\end{enumerate}

\textbf{Questions for the team about this list:}

\begin{enumerate}[resume]
\item Am I right that \#34 and \#35 are the critical path? Or should
I be running \#57 (multi-threaded benchmark) first, since it doesn't
need vLLM integration and it's the cleanest demonstration of the
dual-QP mechanism?

\item Should I be doing \#37 (24-hour stability) earlier? It's cheap
to set up and it catches bugs that block every other benchmark.

\item Is there an issue missing from this list that you'd add?

\item If I get stuck on \#34 (the vLLM integration), what's the
fallback path? Can I ship a meaningful result with just \#57 +
\#19 + \#58?
\end{enumerate}

\section{Things I don't know I don't know}

\begin{enumerate}
\item What's the single question you wish an MS student in this
space had asked you two years earlier?

\item Among the commercial/open-source KV-transport projects I should
be tracking, which one has the closest failure mode to what I'm
building?

\item Is there a specific benchmark result (not a paper, a JSON file)
that I should be reproducing before I trust my own numbers?

\item What's the most common way a thesis like this fails at defense,
and how do I immunize against it now rather than in nine months?
\end{enumerate}

\vspace{8pt}
\hrule
\vspace{6pt}

\section{Bottleneck inventory (top to bottom)}

Every layer in the disaggregated inference stack where something can
stall, starve, drop, or serialize. Ordered from the user-visible
request down to the NIC wire. Each item is where I expect to see a
problem --- I want to know which ones actually bite in production
and which ones are noise.

\subsection*{Layer 0 --- Request arrival}
\begin{enumerate}
\item \textbf{Request rate variance.} Bursty arrivals collapse
continuous batching into serial execution. What's the production
arrival shape that breaks vLLM scheduling assumptions?
\item \textbf{Prompt length distribution.} A 128-token prompt and a
32k-token prompt produce KV caches that differ by 250$\times$. Does
the WFA classifier need to see prompt length, not just the per-block
size?
\item \textbf{Session continuity.} Multi-turn conversations want
prefix-cache hits; random-access RAG queries don't. How does the
scheduler know which shape it's in?
\end{enumerate}

\subsection*{Layer 1 --- vLLM scheduler and dispatch}
\begin{enumerate}[resume]
\item \textbf{Scheduler tick latency.} vLLM makes placement decisions
every step. What does that latency look like on our boxes, and does
it bound the RDMA dispatch budget from above?
\item \textbf{Preemption of long decodes.} When a high-priority
request arrives mid-decode, the scheduler can preempt. Does that
preemption flush KV state we'd otherwise want to reuse?
\item \textbf{Batch shape re-decision.} Prefix cache hits change the
compute shape of a batch. Does the scheduler see enough to adapt, or
does it ride the original batch plan?
\end{enumerate}

\subsection*{Layer 2 --- Model compute on GPU}
\begin{enumerate}[resume]
\item \textbf{Attention compute ($O(n^2 d)$).} At 128k context, a
single attention layer dominates the decode step. Is the transport
layer invisible until attention is already solved by FlashAttention?
\item \textbf{FFN matmul.} For dense models, FFN is usually compute-bound;
for MoE it becomes expert-shuffle-bound. Which regime is each model
in on our hardware?
\item \textbf{Tensor-core utilization.} Ampere (3090) uses FP16/BF16
tensor cores; Blackwell (5090) adds FP8/MXFP4. Is our transport
budget sized against the slower GPU or the faster one?
\item \textbf{Kernel launch overhead.} For small batches, launch
overhead can exceed compute. Is that why my 4-layer Granite bench
shows more variance than the 40-layer one?
\end{enumerate}

\subsection*{Layer 3 --- KV cache lifecycle}
\begin{enumerate}[resume]
\item \textbf{KV cache growth.} Linear in context length. At 128k
context on a 7B model, $\approx 16$\,GB per user. Per-user, not
aggregate.
\item \textbf{PagedAttention block allocation.} 16 tokens per block
at default vLLM settings; the block table fragments as requests
enter and exit.
\item \textbf{Block eviction / swap.} When HBM fills, blocks spill
to CPU RAM or disk. Which blocks get evicted first under the current
policy, and does the policy match what attention actually reads?
\item \textbf{Prefix-cache hit rate.} Cross-request reuse of
already-computed KV blocks. High hit rate flips the workload from
``compute-bound'' to ``cache-transport-bound.''
\item \textbf{Cross-node KV transfer.} The core thesis problem. A
16\,MB block on 10\,GbE takes $\approx 13$\,ms at line rate --- an
order of magnitude longer than the hot-path dispatch budget.
\end{enumerate}

\subsection*{Layer 4 --- Memory blocks and HBM}
\begin{enumerate}[resume]
\item \textbf{HBM capacity.} 24\,GB per 3090, 32\,GB per 5090.
Model weights + activations + KV eat this fast at long context.
\item \textbf{HBM bandwidth ($\approx$ 800 GB/s).} Setting the
ceiling for on-GPU data movement. FlashAttention's tiling exists to
stay inside this budget.
\item \textbf{HBM fragmentation.} Repeated alloc/free of
variably-sized KV blocks fragments the GPU allocator. PagedAttention
mitigates via fixed block sizes, but the fragmentation shows up at
the block-table level.
\item \textbf{Host RAM pinned pages.} RDMA MRs require pinned
memory. The \texttt{ulimit} for locked pages is a hidden ceiling
people hit at production scale.
\item \textbf{RDMA buffer slab fragmentation.} A slab registered
once, sliced into blocks: how do reclamation holes behave under
long-running load?
\end{enumerate}

\subsection*{Layer 5 --- Tensor access and gather/scatter}
\begin{enumerate}[resume]
\item \textbf{Coalesced vs strided reads.} A per-layer KV transfer
touches heads $\times$ seq $\times$ head\_dim. Is that a single
contiguous slab on GPU, or does the layout force strided access?
\item \textbf{Scatter-gather for paged KV.} PagedAttention blocks
are not physically contiguous. RDMA WRITE with an SGE list vs one
WRITE per block --- where's the crossover?
\item \textbf{GPU $\rightarrow$ host memcpy.} Without GDR we pay a
PCIe round-trip to stage KV through host RAM. What fraction of the
hot-path 8\,$\mu$s is actually the memcpy?
\item \textbf{Attention gather pattern.} Sparse attention (sliding
window, sinks, MoE expert routing) gathers an irregular subset of
tokens. Does irregular gather compose with RDMA WRITE, or does it
force SEND/RECV?
\end{enumerate}

\subsection*{Layer 6 --- Host CPU and PCIe}
\begin{enumerate}[resume]
\item \textbf{Python GIL.} pyverbs posts work under the GIL. At a
few hundred posts per second is that measurable?
\item \textbf{CPU core pinning.} Busy-poll CQ thread steals a core.
On a 3090 box serving 3 decode workers, that's a 25\,\% penalty
before any RDMA work happens.
\item \textbf{NUMA crossing.} NIC and GPU on different sockets = PCIe
peer-to-peer forced to go through the CPU. Consumer boards don't
always give us control over this.
\item \textbf{IRQ affinity.} NIC IRQs on the wrong core serialize
completion handling with application work.
\item \textbf{PCIe bus contention.} GPU and NIC share PCIe root
complexes. High GPU I/O + high NIC I/O fight for lanes.
\item \textbf{GPUDirect RDMA absence.} Without GDR (we don't have
BlueField), every KV block stages through host RAM. Latency floor
rises by the PCIe round-trip.
\end{enumerate}

\subsection*{Layer 7 --- RDMA memory registration}
\begin{enumerate}[resume]
\item \textbf{MR registration cost.} First registration is slow; MR
cache avoids repeating it. But cache invalidation on realloc is
expensive.
\item \textbf{ODP (on-demand paging).} Avoids pre-registration but
pays per-page-fault latency. Is it worth it for variable-length
KV buffers?
\item \textbf{rkey rotation.} A long-lived MR leaks its rkey to every
peer. Periodic re-registration closes that hole but costs cycles.
\end{enumerate}

\subsection*{Layer 8 --- RDMA transport (the thesis layer)}
\begin{enumerate}[resume]
\item \textbf{Single-QP head-of-line blocking.} The problem the
dual-QP pool targets. An 8\,$\mu$s metadata op stuck behind a 13\,ms
KV block serialization.
\item \textbf{Send queue depth.} Too shallow: starvation. Too deep:
queuing delay defeats the hot/cold split.
\item \textbf{CQ depth.} Overflow = dropped completions = undefined
behavior. What's a safe multiple of SQ depth?
\item \textbf{Completion polling mode.} Busy-poll burns CPU;
event-driven adds wake-up latency. Adaptive policies exist but are
tricky to tune.
\item \textbf{Signaled-every-N completions.} Unsignaled completions
reduce CQ pressure but delay error detection.
\item \textbf{Inline vs non-inline sends.} Below the inline limit
(64\,B typical), the payload rides the WR; above it, a DMA is
required.
\item \textbf{WR batching.} Posting WRs in a linked list amortizes
doorbell writes. At what batch size does it help?
\item \textbf{RDMA READ vs WRITE for KV fetch.} READ is receiver-paced;
WRITE is sender-paced. Wrong choice inverts the flow control I want.
\item \textbf{Credit-based flow control (RC).} RC QPs have an implicit
credit window. Deep send queues against a slow receiver exhaust
credits silently.
\item \textbf{Ordering guarantees.} RC preserves per-QP order; the
dual-QP pool breaks cross-QP ordering. Does any upstream consumer
depend on it?
\end{enumerate}

\subsection*{Layer 9 --- NIC hardware and PCIe}
\begin{enumerate}[resume]
\item \textbf{NIC work-request rate ceiling.} There's a per-device
WRs/sec number that hardware can sustain. For consumer NICs
(AQC107, X710) I don't know what that is.
\item \textbf{NIC TX/RX ring depth.} Small rings drop under bursty
load before the CQ has a chance to notice.
\item \textbf{MTU and segmentation.} Large KV blocks are segmented
at MTU (1500 default, 4096 or 9000 jumbo). Every segment boundary is
a wire-latency increment.
\item \textbf{PCIe $\leftrightarrow$ NIC bottleneck.} A 10\,GbE NIC
on a PCIe Gen2 x4 link is fine; on Gen4 x8 it's over-provisioned.
Worth checking what my boards actually wire.
\item \textbf{Port counters.} \texttt{perfquery} surfaces link-layer
errors silently accumulating. Worth polling continuously.
\end{enumerate}

\subsection*{Layer 10 --- Wire and link layer}
\begin{enumerate}[resume]
\item \textbf{10\,GbE line rate ($\approx$ 1.25 GB/s).} Hard ceiling
for cold-path throughput on our cluster. No way around it without
hardware change.
\item \textbf{Packet drops.} A single drop triggers a retransmit that
dwarfs the hot-path budget. Direct-connect makes drops rare but not
zero.
\item \textbf{PFC / ECN / DCQCN.} Production RoCE clusters rely on
PFC for lossless behavior. Without a switch we lose the PFC story.
\item \textbf{Switched vs direct-connect.} Direct-connect eliminates
switch hop latency ($\approx$ 1\,$\mu$s per hop) but introduces the
ARP/AH wireup problem that PR~\#11307 addresses.
\item \textbf{Cable quality / CRC errors.} Cheap 10GBASE-T cables
silently corrupt under load. Enterprise clusters filter for this;
mine doesn't.
\end{enumerate}

\subsection*{Layer 11 --- Wireup and connection management}
\begin{enumerate}[resume]
\item \textbf{TCP bootstrap handshake.} Fixed-size 64-byte QpInfo
exchange with retry. My port of libmesh-rdma's protocol; the
security audit hardening lives here.
\item \textbf{GID discovery.} IPv4-mapped GID scan. Hard-coded
index 0 is fragile; we scan the GID table on startup.
\item \textbf{QP state transitions.} RESET $\rightarrow$ INIT
$\rightarrow$ RTR $\rightarrow$ RTS with exponential-backoff retry.
Transient failures show up here first.
\item \textbf{Address resolution (ARP/AH).} The r/HPC bug: AH
creation needs L2 resolution, which needs a shared subnet or a
switch. RC QPs skip this by using the remote GID directly --- the
observation that motivated PR~\#11307.
\end{enumerate}

\subsection*{Layer 12 --- Compression (scuffedQuant)}
\begin{enumerate}[resume]
\item \textbf{CPU Walsh-Hadamard compute.} $O(d \log d)$ per vector.
At 40 layers $\times$ 8 heads $\times$ seq $\times$ 64 head\_dim this
adds up fast on the critical path.
\item \textbf{Lloyd-Max codebook lookup.} Binary search over
$2^b$ levels. Small constant; rarely the bottleneck.
\item \textbf{QJL sketch generation.} Random-projection $\times$
sign-extraction. Quadratic-ish in sketch dimension.
\item \textbf{GPU decompression.} We run this on the decode side to
reconstruct FP32. Matmul-based Hadamard is bandwidth-bound on the
codebook gather, not compute-bound.
\item \textbf{Compression in the critical path.} Every 3-bit
compression ratio we win adds CPU cycles. At what compression savings
does the CPU cost stop being free?
\end{enumerate}

\subsection*{Layer 13 --- Classifier and control plane}
\begin{enumerate}[resume]
\item \textbf{WFA classifier latency.} Per-transfer classification
adds work before every \texttt{ibv\_post\_send}. Microseconds at
most, but measurable on the hot path.
\item \textbf{PMP control loop frequency.} The bang-bang switching
decision runs per transfer. A 1\,$\mu$s loop on every hot-QP post is
an 8\,\% overhead against the 8\,$\mu$s baseline.
\item \textbf{Metric freshness.} Queue depth and service rate feed
PMP. If readings lag by milliseconds, PMP optimizes for yesterday's
load.
\item \textbf{Priority inversion.} A low-priority cold transfer
holding an MR that a high-priority hot transfer wants. Classic
priority-inversion bug; does the dual-QP pool actually prevent it?
\end{enumerate}

\subsection*{Layer 14 --- Observability and logging}
\begin{enumerate}[resume]
\item \textbf{Metric collection overhead.} Scraping
\texttt{ibv\_query\_qp} and \texttt{nvidia-smi} at high frequency
steals cycles from the hot path.
\item \textbf{Structured logging I/O.} Every classified transfer logs
a decision. Disk I/O backpressure is a silent killer of latency
tails.
\item \textbf{Prometheus-style scrape vs event-driven push.} Scrape
interval defines metric staleness. At 15\,s intervals, the scheduler
reacts to load from 15\,s ago.
\end{enumerate}

\subsection*{The three bottlenecks I expect to matter most}

If I had to rank: (a) single-QP head-of-line blocking (Layer~8) ---
the thesis thesis, (b) cross-node KV transfer time (Layer~3), and
(c) PCIe round-trip without GPUDirect (Layer~6). Every other layer
is either solved by someone else (attention, PagedAttention,
FlashAttention), bounded by hardware we can't change (wire rate,
HBM), or small enough to be noise until (a)--(c) are addressed.

\textbf{Question for the team:} is my ranking right? What should
bump up the list that I'm underestimating?

\vspace{8pt}
\hrule
\vspace{6pt}

\section*{Housekeeping}

\begin{itemize}
\item Repository: \url{https://github.com/ndg8743/ScuffedRDMA}
\item Thesis updates are in \texttt{Updates/} --- read Update~4 and
Update~5 before a deep-dive call.
\item All measurements in the \textbf{Results} table in the README have
a JSON file next to them. Anything else I quote is a target, not a
measurement.
\item My advisors: Dr.\ Wafi Danesh, Dr.\ Ashley Suchy, SUNY New
Paltz.
\end{itemize}

\end{document}
